% from ega4-2-fr.tex
\oldpage[IV-2]{147}
pourvu que l'on prouve que $p^{-1}(x)$ est un \emph{préschéma régulier}; autrement
dit, on est ramené à prouver le
\end{proof}

\begin{lemma}[6.7.4.1]
\label{IV.2.s6.7.4.1-fr}
Soient $K$, $L$ deux extensions d'un corps $k$, dont l'une est de type
fini. Si $L$ est une extension séparable de $k$, l'anneau $K\otimes_kL$
est régulier.
\end{lemma}

\begin{proof}
Distinguons deux cas:

\begin{enumerate}
\item[A)] $L$ est une extension \emph{de type fini} de $k$. Alors, avec les
notations de (6.7.1.1), on peut supposer que $L$ est une extension \emph{finie
séparable} de $k(\mathbf t)$. Pour tout $s\in\operatorname{Spec}(A)$,
$\boldsymbol{k}(s)\otimes_{k(\mathbf t)}L$ est alors composé direct d'un nombre
fini de corps, extensions finies séparables de $\boldsymbol{k}(s)$, donc un anneau
\emph{régulier}; il résulte alors de (6.5.2, (ii)) qu'il suffit de prouver que
l'anneau $A$ est régulier; comme $A=S^{-1}A'$, il suffit de prouver que
$A'$ est régulier (\textbf{0}, 17.3.6), mais on a vu qu'il en était bien ainsi dans
la démonstration de (6.7.1.1).
\item[B)] Soit $(L_\lambda)$ la famille filtrante des sous-extensions de
$L$ qui sont de type fini sur $k$; en vertu de A), chacun des anneaux
$C_\lambda=K\otimes_kL_\lambda$ est régulier; d'autre part, pour
$L_\lambda\subset L_\mu$, $L_\mu$ est un $L_\lambda$-module plat, donc
$C_\mu$ est un $C_\lambda$-module plat $(\mathbf{0}_{\mathrm I},6.2.1)$. Comme
$C=K\otimes_kL=\varinjlim(K\otimes_kL_\lambda)$ est noethérien, on peut
appliquer (5.13.7), et $C$ est donc régulier.
\end{enumerate}
\end{proof}

\begin{remarks}[6.7.5]
\label{IV.2.s6.7.5-fr}
\begin{enumerate}
\item[(i)] Dans la démonstration du fait que $\boldsymbol{P}(X,x)$ entraîne
$\boldsymbol{P}(X',x')$, on ne peut se dispenser de l'hypothèse que $k'$ est une
extension \emph{séparable} de $k$. On l'a déjà vu (4.6.1) lorsque $\boldsymbol{P}$ est la
propriété \emph{d)}; mais même lorsque $X$ est \emph{géométriquement intègre} (4.6.2),
il peut se faire que $X$ soit \emph{régulier} sans que $X'$ soit \emph{normal}.

Prenons par exemple pour $X$ une courbe algébrique normale sur $k$
(\textbf{II}, 7.4.2); les anneaux locaux de $X$ étant intégralement clos et de
dimension $1$ sont des anneaux de valuation discrète, donc réguliers
(\textbf{II}, 7.1.6), et $X$ est donc un $k$-schéma régulier (et \emph{a fortiori}
vérifie $(R_n)$ pour tout $n\geqslant0$). Dire que $X$ est géométriquement
intègre signifie alors que le corps $K$ des fonctions rationnelles sur
$X$ est une extension séparable et primaire de $k$ (4.6.3).
Or, prenons pour $k$ un corps non parfait de caractéristique $p>2$, et
soit $a\in k$ un élément n'appartenant pas à $k^p$. Soit
$B$ l'anneau de polynômes $k[S,T]$ en deux indéterminées $S$, $T$; le
polynôme $P(S,T)=T^2\setmin S^p+a$
est irréductible dans $B$, car on vérifie aussitôt que $S^p\setmin a$ n'est pas
un carré dans l'anneau $k[S]$; donc $X=\operatorname{Spec}(A)$, où $A=B/PB$,
est une courbe affine irréductible sur $k$. Pour voir que le schéma $X$
est régulier, il suffit de montrer qu'il est \emph{normal} (\textbf{II}, 7.4.5); or
$A=k[S][t]$, où $t$ est une racine du polynôme $P$ considéré comme
polynôme en $T$ sur $k[S]$, de sorte que le corps $K=R(X)$ des fonctions
rationnelles sur $X$ est le corps des fractions $k(S)[T]/(P)$ de $A$,
extension quadratique de $k(S)$, donc \emph{séparable} sur $k(S)$, et \emph{a fortiori}
sur $k$. Comme $2$ est inversible dans $k$, on vérifie aussitôt que $A$
est la fermeture intégrale de $k[S]$ dans $K$, donc est intégralement clos,
ce qui montre que $X$ est \emph{régulier}. Par ailleurs, si un élément $f+tg$
de $K$ (avec $f$, $g$ dans $k(S)$) est algébrique sur $k$, il en est de même de
sa norme et de sa trace sur $k(S)$, et comme $k(S)$ est une extension pure
de $k$, on en conclut aisément que l'on doit avoir $f\in k$ et $g=0$,
autrement dit $k$ est \emph{algébriquement fermé} dans $K$, et \emph{a fortiori} $K$
est une extension \emph{primaire}
\oldpage[IV-2]{148}
de $k$ (4.3.1). Cependant, si
$k'=k(a^{1/p})$, $X'=X\otimes_k k'$ n'est pas normal, car dans $k'[S]$
on peut écrire $S^p\setmin a=(S\setmin a^{1/p})^p$, et $X'$ est donc isomorphe à
$\operatorname{Spec}(A')$, où $A'=k'[S][t']$, $t'$ étant racine du
polynôme $T^2\setmin S^p$. Or, $A'$ n'est pas intégralement clos, car l'élément
$t''=t'/S$ du corps des fractions $K'=k'(S)[t']$ de $A'$ vérifie
l'équation de dépendance intégrale $t''^2\setmin S^{p-2}=0$ sur $A'$ et
n'appartient pas à $A'$. Dans la théorie classique, cela s'exprime en
disant que le «\emph{genre}» sur $k'$ du corps des fonctions rationnelles $K'$ de
$X'$ est \emph{strictement inférieur} à celui de $K$ sur $k$.

